\documentclass[10pt,a4paper]{amsart}
\usepackage[margin=19mm]{geometry}
\usepackage{amsmath,amssymb,mathtools}
\usepackage[scr=rsfs,cal=euler]{mathalfa}
\usepackage{xfrac}
\usepackage[unicode,hidelinks,bookmarks=false]{hyperref}
\newcommand{\ie}{i.e.\ }
\newcommand{\cf}{cf.\ }
\newcommand{\resp}{respectively\ }
\newcommand{\Subsection}{\subsection}
\providecommand{\nobreakdash}{\nobreak\hbox{-}}
\setlength{\emergencystretch}{2em}
\multlinegap0pt
\usepackage{bm}
\usepackage{bbm}
\usepackage{dsfont}
\usepackage{xspace}
\usepackage{cancel}
\usepackage{mathtools}
\usepackage{amsthm}
\makeatletter
\@ifundefined{theorem}{\newtheorem{theorem}{Theorem}}{}
\@ifundefined{proposition}{\newtheorem{proposition}[theorem]{Proposition}}{}
\makeatother
\title[On two-dimensional steady compactly supported Euler flows wi]{On two-dimensional steady compactly supported Euler flows with constant vorticity}
\author{Changfeng Gui, Jun Wang, Wen Yang, Yong Zhang}
\date{Source: arXiv:2602.07407v4 (CC BY 4.0)}
\begin{document}
\maketitle

\noindent\emph{Source and licence.} On two-dimensional steady compactly supported Euler flows with constant vorticity. Version arXiv:2602.07407v4 (CC BY 4.0), distributed under the Creative Commons Attribution 4.0 International licence, as recorded in the arXiv licence metadata for this exact version and shown on the abstract page https://arxiv.org/abs/2602.07407v4. Full rights notice: licenses/arxiv-2602.07407-CC-BY-4.0.txt.\par\smallskip

\noindent\textit{Editorial scope.} Counted unit: the Proposition whose source label is \texttt{pro-5.3} in the article (source0.tex lines 1276--1287, source label \texttt{pro-5.3}), delivered together with the article's own complete original proof of it (source0.tex lines 1289--1328, one \texttt{proof} environment of 40 lines). No mathematics is written, repaired or re-derived: statement and proof are the article's own text line for line. Required context retained uncounted: e3.5 (lines 652--658); pro-3.3 (lines 673--679); e5.3 (lines 1212--1217); e5.4 (lines 1223--1225); e5.5 (lines 1256--1267). Every definition, display and label of those lines is retained unchanged. Omitted from this package and not redistributed: 1--651, 659--672, 680--1211, 1218--1222, 1226--1255, 1268--1275, 1288--1288, 1329--1741. Nothing inside a retained excerpt is omitted or altered: every retained body is delivered exactly as the article's lines are written, so no omission receipt of any kind accompanies this package. Citations: the retained excerpts carry no citation command at all, so no bibliography entry of the article is transcribed and no reference is redistributed. Reference adaptation: none was needed and none was made; every reference inside the delivered unit resolves inside it, so no unresolved reference or citation is produced. Printed slips kept literal: no printed word, symbol, hypothesis or number of the counted statement or of its proof was altered. Layout and class: the article's own class is \texttt{amsart} with packages including amssymb, latexsym, amsmath, color, amsthm, mathrsfs, srcltx, epsfig, graphicx, hyperref, titletoc, refcheck, palatino, mathpazo; the wrapper is the lane's pinned \texttt{amsart} editorial wrapper at 10pt on A4, which restates the theorem-like environments the retained text needs and adds the article's own macro definitions that the retained text needs (none), with extra wrapper packages (bm, bbm, dsfont, xspace, cancel, mathtools). The delivered numbering is the unit's own. Assets: no retained span contains a figure environment, a graphics-inclusion command, a PDF-page inclusion, a table or a TikZ picture; no image file is copied into this package, the package declares no asset dependency and no image file is redistributed with it. Licence: version arXiv:2602.07407v4 of this article is distributed under the Creative Commons Attribution 4.0 International licence, as recorded in the arXiv licence metadata for this exact version and shown on the abstract page https://arxiv.org/abs/2602.07407v4. A full rights notice is delivered at licenses/arxiv-2602.07407-CC-BY-4.0.txt. Characters: every retained excerpt is pure ASCII, so the delivered engine, XeLaTeX with its default Latin Modern text font, reports no missing character.
\begin{equation}\label{e3.5}
\begin{cases}
\Delta \psi' = 0 & \text{in } \Omega_0 \setminus \overline{B_\lambda}, \\
\psi' = -\partial_\nu \psi^\gamma_\mathrm{tri} \, \eta_0 = -\sqrt{Q^\gamma} \, \eta_0 & \text{on } \partial \Omega_0, \\
\psi' = 0 & \text{on } \partial B_\lambda.
\end{cases}
\end{equation}

\begin{proposition}\label{pro-3.3}
The map \(G\) is Fr\'{e}chet differentiable in a neighborhood of \((\gamma, 0)\). Moreover,
\begin{equation}\label{e3.6}
\partial_\eta G(\gamma, 0)[\eta_0] = 2 \partial_\nu \psi^\gamma_\mathrm{tri} \bigl( \partial_\nu \psi' + \partial_{\nu\nu} \psi^\gamma_\mathrm{tri} \, \eta_0 \bigr), \quad \eta_0 \in X,
\end{equation}
where \(\psi'\) solves \eqref{e3.5}.
\end{proposition}

\begin{equation}\label{e5.3}
\begin{cases}
\widetilde{\psi}_{1triv}(|x|) = -\frac{\left(1-|x|^2\right)\widetilde{\gamma}}{4}, & |x| \in [0,1), \\
\widetilde{\psi}_{2triv}(|x|) = a\ln |x|, & |x| \in [1,+\infty),
\end{cases}
\end{equation}

\begin{equation}\label{e5.4}
\widetilde{F}(\widetilde{\gamma},\eta) =\Pi_0\left( \big| \partial_{\nu_\eta} \widetilde{\psi}_{1\eta}\big|_{\partial \Omega_0}\big|^2+\beta\mathcal{K}_{\eta}-q\big| \widetilde{\partial}_{\nu_\eta} \psi_{2\eta}\big|_{\partial \Omega_0}\big|^2-Q(\widetilde{\gamma})\right) ,
\end{equation}

\begin{equation}\label{e5.5}
\begin{cases}
\Delta \widetilde{\psi}'_{1} = 0 & \text{in } \Omega_0, \\
\widetilde{\psi}'_{1} = -\partial_{\nu}\widetilde{\psi}_{1triv} \eta_0 & \text{on } \partial \Omega_0,
\end{cases}
\qquad
\begin{cases}
\Delta \widetilde{\psi}'_{2} = 0 & \text{in } \mathbb{R}^2 \setminus \overline{\Omega_0}, \\
\widetilde{\psi}'_{2} = -\partial_{\nu}\widetilde{\psi}_{2triv} \eta_0 & \text{on } \partial \Omega_0, \\
\nabla\widetilde{\psi}'_{2} \rightarrow0 & \text{as } |x|\rightarrow +\infty,
\end{cases}
\end{equation}

\begin{proposition}\label{pro-5.3}
The map \(\widetilde{F}\) $($defined in \eqref{e5.4}$)$ is Fr\'{e}chet differentiable near \((\beta,0)\). Moreover, for any \(\eta_0 \in \widetilde{X}\), there holds that
\begin{equation}\label{e5.6}
\begin{aligned}
\partial_{\eta} \widetilde{F}(\beta,0)[\eta_0] =&2\partial_{\nu}\widetilde{\psi}_{1triv}\left(\partial_{\nu}\widetilde{\psi}'_{1}+\partial_{\nu\nu}\widetilde{\psi}_{1triv} \eta_0\right)
-\beta(\partial_{\theta\theta}+1)\eta_0\\
-&2q\partial_{\nu}\widetilde{\psi}_{2triv}\left(\partial_{\nu}\widetilde{\psi}'_{2}+
\partial_{\nu\nu}\widetilde{\psi}_{2triv} \eta_0\right),
\end{aligned}
\end{equation}
where $\widetilde{\psi}_{1triv}$ and $\widetilde{\psi}_{2triv}$ are trivial solutions given in \eqref{e5.3} and \(\widetilde{\psi}'_{1}\) and \(\widetilde{\psi}'_{2}\) solve \eqref{e5.5}, respectively.
\end{proposition}

\begin{proof}
We compute the G\^{a}teaux derivative with respect to \(\eta\):
\[
\partial_\eta \widetilde{F}(\beta,0)[\eta_0] = \frac{d}{dt}\Big|_{t=0} \widetilde{F}(\beta,t\eta_0).
\]
Following the same computational procedure as in the proof of Proposition \ref{pro-3.3}, the first and third terms on the right-hand side of \eqref{e5.6} are readily obtained. We now focus mainly on the linearity of the curvature term. Since the curvature of a planar curve represented as a radial graph is given by
\[
\mathcal{K}_{\eta}
=
\frac{(1+\eta)^2+2(\partial_{\theta}\eta)^2
-(1+\eta)\partial_{\theta\theta}\eta}
{\left((1+\eta)^2+(\partial_{\theta}\eta)^2\right)^{3/2}},
\]
we may expand it around the unit disk. Indeed, if the perturbation $\eta$ is sufficiently small, then
\[
\mathcal{K}_{\eta}
=
\frac{
1+2\eta-\partial_{\theta\theta}\eta
+O\!\left(\eta^2+(\partial_{\theta}\eta)^2\right)
}
{
1+3\eta
+O\!\left(\eta^2+(\partial_{\theta}\eta)^2\right)
},
\]
which yields
\[
\mathcal{K}_{\eta}
=
1-\eta-\partial_{\theta\theta}\eta
+O\!\left(\eta^2+(\partial_{\theta}\eta)^2\right).
\]
Therefore, the linearization of the curvature operator at the unit disk is given by
\[
\partial_{\eta}\mathcal{K}(0)[\eta_0]
=
-\left(\partial_{\theta\theta}+1\right)\eta_0.
\]
\end{proof}

\end{document}
